\documentclass{article}
\usepackage[T1]{fontenc}
\usepackage{lmodern}
\usepackage{graphicx}
\usepackage{amsmath,amssymb}
\usepackage[a4paper,margin=2cm]{geometry}
\usepackage[absolute,overlay]{textpos}
\setlength{\TPHorizModule}{1mm}
\setlength{\TPVertModule}{1mm}
\usepackage[indonesian]{babel}
\usepackage{hyperref}
\setlength{\emergencystretch}{2em}
% unit: Halaman 52

\begin{document}

% === Halaman 52 (llm) ===

\section*{Serangan lainnya: Side-Channel Attack}

\begin{itemize}
    \item Serangan ini tidak terutama menyerang algoritma matematisnya, tetapi memanfaatkan informasi yang bocor ketika algoritma dijalankan pada perangkat fisik.
    \item Ada beberapa macam side-channel attack:
    \begin{enumerate}
        \item \textbf{Timing Attack} \\ Memanfaatkan perbedaan waktu eksekusi untuk memperoleh informasi tentang kunci.
        \item \textbf{Power Analysis Attack} \\ Menganalisis konsumsi daya perangkat ketika melakukan operasi kriptografi. Contohnya:
        \begin{itemize}
            \item Simple Power Analysis (SPA)
            \item Differential Power Analysis (DPA)
        \end{itemize}
        \item \textbf{Electromagnetic (EM) Attack} \\ Menganalisis radiasi elektromagnetik yang dihasilkan perangkat ketika melakukan operasi kriptografi.
        \item \textbf{Fault Injection Attack} \\ Penyerang sengaja menyebabkan kesalahan pada perangkat atau proses komputasi, kemudian menganalisis \textit{output} yang salah tersebut. Salah satu contohnya adalah Differential Fault Analysis (DFA).
    \end{enumerate}
\end{itemize}

\end{document}
